% slide-ready summary for the lecture box 'Research from our group' (max 6 lines); generated by the notebook
\begin{itemize}\setlength{\itemsep}{2pt}
  \item \textbf{Elasticitatea cererii de energie electrică în România} (Proșcanu, Proșcanu și Pele, 2025, JSES): elasticitate de 0,056 (mică și pozitivă; Tabelul~7); cauzalitatea Granger merge de la consum la preț, nu invers (Tabelele~4--5).
  \item Replicarea noastră pe aceleași 72 de luni: 0,064; cu prețul gazelor ca instrument, $-$0,023 --- cererea aproape nu reacționează la preț.
  \item Decizia: nu există o acoperire naturală --- un furnizor cu preț fix ar trebui să cumpere la termen $\approx$0,93 din volumul așteptat (riscul scade cu 92\,\% în eșantion; în afara eșantionului, $h^*$ = 1,06 estimat pe 2016--2020 îl reduce cu 81\,\% în 2021--2025).
\end{itemize}
